L'albédo de la Terre est d'environ \num{\v{a1}}. La fonte des glaces et de la neige, très réfléchissantes, remplace des surfaces claires par de l'eau ou du sol plus sombres. Supposons que cela fasse baisser l'albédo terrestre de \num{\v{da}}. La constante solaire est de \qty{1361}{W/m^2}.

\begin{enumerate}[label=\alph*)]
    \item Déterminez la température d'équilibre de la Terre (sans effet de serre) avec l'albédo initial.
    \item Déterminez la nouvelle température d'équilibre.
    \item De combien la température d'équilibre augmente-t-elle ? Expliquez pourquoi on parle d'une « rétroaction positive ».
\end{enumerate}
